\chapter{Phase 1: Preparing the Raspberry Pi}

\section{Goal}

By the end of this phase, the Raspberry Pi must be:

\begin{itemize}
  \item Booting Raspberry Pi OS Lite (64-bit).
  \item Reachable over SSH from the working Mac.
  \item Assigned a static IP reserved on the router.
  \item Running Asterisk 22 and responding on its CLI (\texttt{chan\_pjsip} module).
  \item With the project's configuration files placed in \texttt{/etc/asterisk/}.
\end{itemize}

\section{Materials}

\begin{itemize}
  \item Raspberry Pi 4 (2\,GB)
  \item 32\,GB microSD
  \item Official USB-C power supply
  \item micro-HDMI cable
  \item USB keyboard (only for the first boot)
  \item Television or HDMI monitor
  \item Mac (or other computer) with microSD reader
\end{itemize}

\section{Step 1: Flash the microSD}

From the Mac, download and install the \emph{Raspberry Pi Imager} from the official site \texttt{raspberrypi.com/software}. Connect the microSD to the Mac (with a USB adapter if necessary).

When the Imager opens, configure:

\begin{itemize}
  \item \textbf{Device:} Raspberry Pi 4
  \item \textbf{Operating System:} Raspberry Pi OS Lite (64-bit)
  \item \textbf{Storage:} the inserted microSD
\end{itemize}

Before clicking \emph{Write}, click the gear icon (\textsf{\textbf{\textcolor{hugoblue}{Edit Settings}}}) to configure:

\begin{itemize}
  \item \textbf{Hostname:} \texttt{hugo-pbx}
  \item \textbf{User:} \texttt{beto} (or whatever name you prefer), with a strong password.
  \item \textbf{SSH:} enabled, password or public-key authentication.
  \item \textbf{WiFi:} SSID and password (only if you won't use Ethernet).
  \item \textbf{Locale:} \texttt{America/Santiago}, \texttt{latam} keyboard.
\end{itemize}

Confirm and let the Imager write (between 5 and 15 minutes depending on the reader's speed).

\section{Step 2: First boot}

Insert the microSD into the Pi, plug in HDMI, the USB keyboard and, finally, the power supply. The first boot takes between 1 and 2 minutes because it automatically resizes the filesystem. When it finishes, a \emph{login prompt} will be shown.

Log in with the username and password configured in the previous step.

\section{Step 3: Network configuration}

The desirable setup is a static IP reserved on the router for the Pi's MAC address. To obtain the MAC of the Ethernet port:

\begin{shellcode}
ip link show eth0 | grep ether
\end{shellcode}

The output will look like \texttt{link/ether dc:a6:32:xx:xx:xx}. Note that MAC, enter the router's admin panel (usually \texttt{http://192.168.1.1}) and reserve the IP \texttt{192.168.1.10} for that MAC.

Reboot the Pi and verify the IP:

\begin{shellcode}
ip addr show eth0
\end{shellcode}

\begin{tip}
If the router hands out IPs from another range (e.g. \texttt{192.168.0.x} or \texttt{10.0.0.x}), adapt every address in the manual accordingly.
\end{tip}

\section{Step 4: SSH connection}

From the Mac:

\begin{shellcode}
ssh beto@192.168.1.10
\end{shellcode}

If the router resolves mDNS, this also works:

\begin{shellcode}
ssh beto@hugo-pbx.local
\end{shellcode}

From this point on the USB keyboard and HDMI cable are no longer needed.

\section{Step 5: System update}

\begin{shellcode}
sudo apt update
sudo apt upgrade -y
sudo reboot
\end{shellcode}

Reconnect over SSH after the reboot.

\section{Step 6: Clone the project repository}

Copy the \texttt{hugo-phone/} directory delivered in the project repository onto the Pi. If it's in a Git service:

\begin{shellcode}
cd ~
git clone <repository-URL> hugo-phone
cd hugo-phone
\end{shellcode}

If the files are only on the Mac:

\begin{shellcode}
# From the Mac
scp -r hugo-phone/ beto@hugo-pbx.local:~/
\end{shellcode}

\section{Step 7: Run the installer}

\begin{shellcode}
cd ~/hugo-phone
sudo bash scripts/install-all.sh
\end{shellcode}

The script performs, in order:

\begin{enumerate}
  \item Full \texttt{apt update \&\& apt upgrade}.
  \item Installation of Asterisk 22, DOSBox, \texttt{build-essential}, \texttt{libudev-dev} and other dependencies.
  \item Copies the \texttt{pjsip.conf}, \texttt{extensions.conf} and \texttt{manager.conf} files to \texttt{/etc/asterisk/} with the correct permissions.
  \item Installation of \emph{rustup} for the \texttt{beto} user (downloads and installs the stable toolchain).
  \item Compilation of the bridge in \texttt{release} mode, with LTO and \texttt{strip}.
  \item Creation of the \texttt{udev} rule for \texttt{/dev/uinput} and loading of the kernel module.
  \item Installation of the \texttt{systemd} service \texttt{hugo-bridge.service}.
  \item Deployment of the DOSBox configuration to \texttt{/opt/hugo-phone/games/dosbox-hugo.conf} (invoked with \texttt{dosbox -conf}).
\end{enumerate}

The slowest phase is compiling the bridge, which on a Pi 4 takes between 5 and 8 minutes. Estimated total time: 20--30 minutes depending on network speed.

\section{Step 8: Configure secrets}

The initial passwords are placeholders and must be replaced before any real use:

\begin{shellcode}
sudo nano /etc/asterisk/pjsip.conf
sudo nano /etc/asterisk/manager.conf
sudo nano /opt/hugo-phone/bridge/config.toml
\end{shellcode}

Replace every occurrence of \texttt{cambiame-101}, \texttt{cambiame-102} and \texttt{cambiame-ami}.

\begin{warning}
The password in \texttt{manager.conf} (section \texttt{[hugo-bridge]}) and the one in \texttt{config.toml} (section \texttt{[ami]}) must be \textbf{identical}. Otherwise the bridge won't be able to authenticate against Asterisk.
\end{warning}

Apply the changes by restarting Asterisk:

\begin{shellcode}
sudo systemctl restart asterisk
\end{shellcode}

\section{Step 9: Verification}

Enter the interactive Asterisk console:

\begin{shellcode}
sudo asterisk -rvvv
\end{shellcode}

Inside the CLI, run:

\begin{plaincode}
hugo-pbx*CLI> pjsip show endpoints
\end{plaincode}

It must list extensions 101 and 102 (the latter is used in the next phase) with status \texttt{Not in use} (no clients registered yet).

\begin{plaincode}
hugo-pbx*CLI> manager show settings
\end{plaincode}

It must show \texttt{Enabled: Yes} and \texttt{Port: 5038}.

\begin{plaincode}
hugo-pbx*CLI> dialplan show hugo-game
\end{plaincode}

It must display extension \texttt{9000} with its \texttt{Answer}, \texttt{Wait}, \texttt{Playback(beep)}, \texttt{Read} and \texttt{UserEvent} steps.

If all three commands respond without errors, \textbf{Phase 1 is complete}.

\section{Troubleshooting}

\subsection{Asterisk does not start}

\begin{shellcode}
sudo journalctl -u asterisk -n 80
\end{shellcode}

It is almost always a syntax error in one of the \texttt{.conf} files, which the log points out with an exact line number.

\subsection{The Pi reboots for no apparent reason or shows random lag}

The power supply is probably insufficient. The Pi shows a yellow lightning bolt icon on screen when it detects \emph{undervoltage}. Switch to the official 5\,V/3\,A power supply.

\subsection{SSH doesn't respond}

\begin{itemize}
  \item Check that the Ethernet cable is plugged in or that WiFi has properly joined the network. Plug in keyboard + HDMI and run \texttt{ip addr}.
  \item Check that the SSH service is running: \texttt{sudo systemctl status ssh}.
  \item Confirm that the router's firewall isn't blocking LAN traffic.
\end{itemize}

\subsection{The installer fails while compiling the bridge}

\begin{itemize}
  \item Retry manually as the \texttt{beto} user:
  \begin{shellcode}
cd /opt/hugo-phone/bridge
~/.cargo/bin/cargo build --release
  \end{shellcode}
  \item If it fails due to lack of memory (unlikely on a Pi 4 2\,GB, but possible), add \emph{swap}:
  \begin{shellcode}
sudo dphys-swapfile swapoff
sudo sed -i 's/CONF_SWAPSIZE=.*/CONF_SWAPSIZE=2048/' \
  /etc/dphys-swapfile
sudo dphys-swapfile setup
sudo dphys-swapfile swapon
  \end{shellcode}
\end{itemize}
